\documentclass[12pt,a4paper]{article}

% ------------------------------------------------------------
% Кодировка, язык и геометрия страницы
% ------------------------------------------------------------
\usepackage[T2A]{fontenc}
\usepackage[utf8]{inputenc}
\usepackage[russian]{babel}
\usepackage{geometry}
\geometry{
    left=25mm,
    right=20mm,
    top=20mm,
    bottom=20mm
}

% ------------------------------------------------------------
% Математика
% ------------------------------------------------------------
\usepackage{amsmath,amssymb,amsthm,mathtools,bm}
\usepackage{booktabs,tabularx,array}
\usepackage{enumitem}

% ------------------------------------------------------------
% Алгоритмы и рисунки
% ------------------------------------------------------------
\usepackage{algorithm}
\usepackage{algpseudocode}
\usepackage{tikz}
\usetikzlibrary{positioning,arrows.meta,calc}

% ------------------------------------------------------------
% Гиперссылки внутри документа
% ------------------------------------------------------------
\usepackage[
    colorlinks=true,
    linkcolor=blue!50!black,
    urlcolor=blue!50!black
]{hyperref}

% ------------------------------------------------------------
% Оформление
% ------------------------------------------------------------
\setlength{\parindent}{1.25cm}
\setlength{\parskip}{0.2em}
\setlist[itemize]{topsep=0.3em,itemsep=0.2em}
\setlist[enumerate]{topsep=0.3em,itemsep=0.2em}

\floatname{algorithm}{Алгоритм}
\algrenewcommand\algorithmicrequire{\textbf{Вход:}}
\algrenewcommand\algorithmicensure{\textbf{Выход:}}
\algrenewcommand\algorithmicfor{\textbf{для}}
\algrenewcommand\algorithmicforall{\textbf{для всех}}
\algrenewcommand\algorithmicwhile{\textbf{пока}}
\algrenewcommand\algorithmicdo{\textbf{выполнять}}
\algrenewcommand\algorithmicif{\textbf{если}}
\algrenewcommand\algorithmicthen{\textbf{то}}
\algrenewcommand\algorithmicelse{\textbf{иначе}}
\algrenewcommand\algorithmicend{\textbf{конец}}
\algrenewcommand\algorithmicreturn{\textbf{вернуть}}

\tikzset{
    block/.style={
        draw,
        rounded corners,
        align=center,
        minimum height=10mm,
        minimum width=24mm,
        fill=blue!5
    },
    fwd/.style={
        -{Latex[length=2mm]},
        thick,
        blue!70!black
    },
    bwd/.style={
        -{Latex[length=2mm]},
        thick,
        dashed,
        red!70!black
    }
}

% ------------------------------------------------------------
% Теоремы
% ------------------------------------------------------------
\theoremstyle{plain}
\newtheorem{theorem}{Теорема}[section]
\newtheorem{lemma}[theorem]{Лемма}
\newtheorem{proposition}[theorem]{Утверждение}
\newtheorem{corollary}[theorem]{Следствие}

\theoremstyle{definition}
\newtheorem{definition}{Определение}[section]
\newtheorem{example}{Пример}[section]

\theoremstyle{remark}
\newtheorem{remark}{Замечание}[section]

% ------------------------------------------------------------
% Сокращения
% ------------------------------------------------------------
\newcommand{\R}{\mathbb{R}}
\newcommand{\E}{\mathbb{E}}
\newcommand{\norm}[1]{\left\lVert #1 \right\rVert}
\newcommand{\abs}[1]{\left| #1 \right|}
\newcommand{\inner}[2]{\left\langle #1,#2 \right\rangle}
\newcommand{\dd}{\,\mathrm{d}}
\newcommand{\argmin}{\operatorname*{arg\,min}}
\newcommand{\softmax}{\operatorname{softmax}}
\newcommand{\diag}{\operatorname{diag}}
\newcommand{\sign}{\operatorname{sign}}

\title{Билет №97 \\ \small Глубокие нейронные сети. Градиентный спуск. Обратное распространение. (ИТМО)}
\author{Сериков Владислав Александрович}
\date{19 июля 2026}

\begin{document}

\maketitle
\tableofcontents
\newpage

% ============================================================
\section{Постановка задачи обучения}
% ============================================================

Пусть дана обучающая выборка
\[
    \mathcal{D}
    =
    \left\{
        \left(x_i,y_i\right)
    \right\}_{i=1}^{N},
\]
где $x_i\in\R^{n_0}$ --- входной объект, а $y_i$ --- требуемый
ответ. Нейронная сеть задаёт параметрическое отображение
\[
    f_\theta\colon \R^{n_0}\to\R^{n_L},
\]
где $\theta$ обозначает совокупность всех обучаемых весов и
смещений.

Для одного объекта определяется функция потерь
\[
    \ell\bigl(f_\theta(x_i),y_i\bigr),
\]
измеряющая отличие предсказания сети от правильного ответа.
Обучение сети обычно формулируется как задача минимизации
эмпирического риска:
\[
    \min_{\theta} F(\theta),
\]
где
\[
    F(\theta)
    =
    \frac{1}{N}
    \sum_{i=1}^{N}
    \ell\bigl(f_\theta(x_i),y_i\bigr)
    +
    \lambda\Omega(\theta).
\]
Здесь:

\begin{itemize}
    \item первая часть функционала измеряет ошибку на обучающей
    выборке;
    \item $\Omega(\theta)$ --- регуляризатор;
    \item $\lambda\geq 0$ --- коэффициент регуляризации.
\end{itemize}

Например, при $L_2$-регуляризации весов
\[
    \Omega(\theta)
    =
    \frac{1}{2}
    \sum_{l=1}^{L}
    \norm{W^{(l)}}_F^2.
\]

Выборку обычно разделяют на три части:

\begin{enumerate}
    \item \textbf{обучающая выборка} используется для изменения
    параметров;
    \item \textbf{валидационная выборка} используется для выбора
    архитектуры и гиперпараметров;
    \item \textbf{тестовая выборка} используется для итоговой
    независимой оценки качества.
\end{enumerate}

Важно различать три компонента обучения:

\begin{itemize}
    \item \textbf{прямое распространение} вычисляет предсказание и
    значение функции потерь;
    \item \textbf{обратное распространение} вычисляет производные
    функции потерь по параметрам;
    \item \textbf{алгоритм оптимизации}, например градиентный спуск,
    использует производные для изменения параметров.
\end{itemize}

% ============================================================
\section{Глубокие нейронные сети}
% ============================================================

\subsection{Искусственный нейрон}

Искусственный нейрон принимает входной вектор
$a\in\R^n$ и вычисляет
\[
    z=w^\top a+b,
    \qquad
    h=\varphi(z),
\]
где

\begin{itemize}
    \item $w\in\R^n$ --- вектор весов;
    \item $b\in\R$ --- смещение;
    \item $z$ --- предактивация;
    \item $\varphi$ --- функция активации;
    \item $h$ --- выход нейрона.
\end{itemize}

Смещение можно включить в вектор весов, дополнив вход единицей:
\[
    \widetilde{a}
    =
    \begin{pmatrix}
        a\\
        1
    \end{pmatrix},
    \qquad
    \widetilde{w}
    =
    \begin{pmatrix}
        w\\
        b
    \end{pmatrix},
    \qquad
    z=\widetilde{w}^{\top}\widetilde{a}.
\]

\subsection{Полносвязный слой}

Пусть слой $l-1$ содержит $n_{l-1}$ выходов, а слой $l$ содержит
$n_l$ нейронов. Тогда
\[
    W^{(l)}\in\R^{n_l\times n_{l-1}},
    \qquad
    b^{(l)}\in\R^{n_l}.
\]

Вычисления слоя имеют вид
\[
    z^{(l)}
    =
    W^{(l)}a^{(l-1)}+b^{(l)},
    \qquad
    a^{(l)}
    =
    \varphi_l\bigl(z^{(l)}\bigr).
\]

Здесь функция активации обычно применяется покомпонентно:
\[
    a_j^{(l)}
    =
    \varphi_l\bigl(z_j^{(l)}\bigr).
\]

\subsection{Многослойный перцептрон}

Многослойный перцептрон, или MLP, представляет собой композицию
слоёв:
\[
\begin{aligned}
    a^{(0)} &= x,\\
    z^{(1)} &= W^{(1)}a^{(0)}+b^{(1)},&
    a^{(1)} &= \varphi_1\bigl(z^{(1)}\bigr),\\
    z^{(2)} &= W^{(2)}a^{(1)}+b^{(2)},&
    a^{(2)} &= \varphi_2\bigl(z^{(2)}\bigr),\\
    &\vdotswithin{=}& &\vdotswithin{=}\\
    z^{(L)} &= W^{(L)}a^{(L-1)}+b^{(L)},&
    a^{(L)} &= \varphi_L\bigl(z^{(L)}\bigr).
\end{aligned}
\]

Выход сети равен
\[
    \widehat{y}=a^{(L)}.
\]

Совокупность параметров:
\[
    \theta
    =
    \left\{
        W^{(1)},b^{(1)},
        \ldots,
        W^{(L)},b^{(L)}
    \right\}.
\]

Число параметров полносвязной сети:
\[
    P
    =
    \sum_{l=1}^{L}
    \left(
        n_l n_{l-1}+n_l
    \right).
\]

Под \textbf{глубиной} обычно понимается число последовательных
обучаемых преобразований, а под \textbf{шириной} --- число нейронов
в слое. Строгой универсальной границы между неглубокими и глубокими
сетями нет.

\begin{center}
\begin{tikzpicture}[node distance=15mm]
    \node[block] (x) {$a^{(0)}=x$};
    \node[block, right=of x] (l1)
    {$z^{(1)}=W^{(1)}a^{(0)}+b^{(1)}$\\
     $a^{(1)}=\varphi_1(z^{(1)})$};
    \node[block, right=of l1] (l2)
    {$\cdots$};
    \node[block, right=of l2] (l3)
    {$z^{(L)}=W^{(L)}a^{(L-1)}+b^{(L)}$\\
     $a^{(L)}=\varphi_L(z^{(L)})$};
    \node[block, right=of l3] (loss)
    {$\ell(a^{(L)},y)$};

    \draw[fwd] (x) -- node[above] {прямой проход} (l1);
    \draw[fwd] (l1) -- (l2);
    \draw[fwd] (l2) -- (l3);
    \draw[fwd] (l3) -- (loss);

    \draw[bwd] (loss) to[bend left=35]
        node[below] {производные} (l3);
    \draw[bwd] (l3) to[bend left=35] (l2);
    \draw[bwd] (l2) to[bend left=35] (l1);
    \draw[bwd] (l1) to[bend left=35] (x);
\end{tikzpicture}
\end{center}

\subsection{Роль нелинейности}

Если все функции активации являются тождественными, то композиция
любого числа слоёв остаётся аффинным отображением. Например,
\[
\begin{aligned}
    W^{(2)}
    \left(
        W^{(1)}x+b^{(1)}
    \right)
    +b^{(2)}
    &=
    \left(
        W^{(2)}W^{(1)}
    \right)x
    +
    \left(
        W^{(2)}b^{(1)}+b^{(2)}
    \right)\\
    &=\widetilde{W}x+\widetilde{b}.
\end{aligned}
\]

Следовательно, без нелинейных функций активации увеличение глубины
не увеличивает класс представимых функций: всю сеть можно заменить
одним аффинным слоем.

\subsection{Основные функции активации}

\begin{center}
\begin{tabularx}{\textwidth}{
    >{\raggedright\arraybackslash}p{0.16\textwidth}
    >{\centering\arraybackslash}p{0.27\textwidth}
    >{\centering\arraybackslash}p{0.25\textwidth}
    X
}
\toprule
\textbf{Название} &
\textbf{Функция} &
\textbf{Производная} &
\textbf{Особенности}\\
\midrule

Линейная &
$\varphi(z)=z$ &
$\varphi'(z)=1$ &
Используется, например, в выходном слое регрессии.\\[1ex]

Сигмоида &
$\displaystyle
\sigma(z)=\frac{1}{1+e^{-z}}
$ &
$\sigma'(z)=\sigma(z)(1-\sigma(z))$ &
Выход лежит в $(0,1)$. При больших $\abs{z}$ производная близка
к нулю.\\[2ex]

Гиперболический тангенс &
$\varphi(z)=\tanh z$ &
$\varphi'(z)=1-\tanh^2z$ &
Выход лежит в $(-1,1)$, функция центрирована относительно нуля,
но также насыщается.\\[1ex]

ReLU &
$\varphi(z)=\max(0,z)$ &
$\varphi'(z)=
\begin{cases}
0,&z<0,\\
1,&z>0
\end{cases}$ &
Проста и не насыщается при $z>0$. В точке $z=0$ производная не
существует; программная реализация выбирает фиксированное значение,
обычно $0$.\\[3ex]

Leaky ReLU &
$\varphi(z)=\max(\alpha z,z)$,
$0<\alpha<1$ &
$\varphi'(z)=
\begin{cases}
\alpha,&z<0,\\
1,&z>0
\end{cases}$ &
Сохраняет ненулевую производную в отрицательной полуплоскости.\\
\bottomrule
\end{tabularx}
\end{center}

\subsection{Выходной слой и функция потерь}

\subsubsection{Регрессия}

Для регрессии часто используют линейный выход:
\[
    \widehat{y}=z^{(L)}
\]
и среднеквадратичную ошибку
\[
    \ell(\widehat{y},y)
    =
    \frac{1}{2}
    \norm{\widehat{y}-y}_2^2.
\]

Её градиент по выходу равен
\[
    \nabla_{\widehat{y}}\ell
    =
    \widehat{y}-y.
\]

\subsubsection{Бинарная классификация}

Пусть последний слой формирует скалярный логит $z$. Вероятность
положительного класса определяется как
\[
    p=\sigma(z).
\]

Бинарная перекрёстная энтропия:
\[
    \ell(p,y)
    =
    -y\log p-(1-y)\log(1-p),
    \qquad y\in\{0,1\}.
\]

Производная потерь непосредственно по логиту:
\[
    \frac{\partial\ell}{\partial z}
    =
    p-y.
\]

\subsubsection{Многоклассовая классификация}

Пусть $z=(z_1,\ldots,z_C)^\top$ --- логиты классов. Преобразование
softmax задаётся формулой
\[
    p_j
    =
    \frac{e^{z_j}}
    {\sum_{r=1}^{C}e^{z_r}}.
\]

Если $y$ --- one-hot-вектор правильного класса, то перекрёстная
энтропия равна
\[
    \ell(p,y)
    =
    -\sum_{j=1}^{C}y_j\log p_j.
\]

Для численной устойчивости следует вычислять
\[
    c=\max_j z_j,
    \qquad
    p_j
    =
    \frac{e^{z_j-c}}
    {\sum_{r=1}^{C}e^{z_r-c}}.
\]
Вычитание одной и той же константы из всех логитов не изменяет
softmax.

\subsection{Основные виды глубоких сетей}

\begin{itemize}
    \item \textbf{Полносвязные сети} применяют матричные аффинные
    преобразования между всеми входами и выходами соседних слоёв.

    \item \textbf{Свёрточные сети} используют локальные фильтры и
    разделение весов между разными пространственными позициями.

    \item \textbf{Рекуррентные сети} многократно применяют один и тот
    же параметризованный переход к элементам последовательности.
    Обратное распространение по развёрнутой рекуррентной сети
    называется обратным распространением во времени.

    \item \textbf{Остаточные сети} используют блоки вида
    \[
        h^{(l+1)}
        =
        h^{(l)}
        +
        F_l\bigl(h^{(l)}\bigr).
    \]

    \item \textbf{Сети с механизмом внимания} вычисляют
    взвешенные комбинации представлений различных элементов
    входной последовательности.
\end{itemize}

Во всех этих случаях вычисление градиента можно представить как
обратный проход по вычислительному графу.

% ============================================================
\section{Производные многомерных функций}
% ============================================================

\subsection{Градиент и производная по направлению}

Пусть $F\colon\R^d\to\R$ дифференцируема. Её градиент:
\[
    \nabla F(x)
    =
    \begin{pmatrix}
        \frac{\partial F}{\partial x_1}\\
        \vdots\\
        \frac{\partial F}{\partial x_d}
    \end{pmatrix}.
\]

Для направления $v\in\R^d$ производная по направлению равна
\[
    D_vF(x)
    =
    \lim_{t\to 0}
    \frac{F(x+tv)-F(x)}{t}
    =
    \nabla F(x)^\top v.
\]

Линейная аппроксимация функции имеет вид
\[
    F(x+\Delta x)
    =
    F(x)
    +
    \nabla F(x)^\top\Delta x
    +
    o\bigl(\norm{\Delta x}\bigr).
\]

\begin{theorem}[Направление наискорейшего локального убывания]
Пусть $F\colon\R^d\to\R$ дифференцируема в точке $x$ и
$\nabla F(x)\neq 0$. Среди всех единичных направлений
$v$, $\norm{v}_2=1$, производная по направлению минимальна при
\[
    v_*
    =
    -\frac{\nabla F(x)}{\norm{\nabla F(x)}_2}.
\]
Минимальное значение производной равно
\[
    D_{v_*}F(x)
    =
    -\norm{\nabla F(x)}_2.
\]
\end{theorem}

\begin{proof}
По формуле производной по направлению
\[
    D_vF(x)
    =
    \inner{\nabla F(x)}{v}.
\]
Из неравенства Коши--Буняковского следует
\[
    \inner{\nabla F(x)}{v}
    \geq
    -\norm{\nabla F(x)}_2\norm{v}_2.
\]
Так как $\norm{v}_2=1$, получаем
\[
    D_vF(x)
    \geq
    -\norm{\nabla F(x)}_2.
\]
Равенство в неравенстве Коши--Буняковского достигается тогда и
только тогда, когда $v$ направлен противоположно градиенту:
\[
    v
    =
    -\frac{\nabla F(x)}
    {\norm{\nabla F(x)}_2}.
\]
Следовательно, отрицательный градиент задаёт направление
наискорейшего локального убывания функции в евклидовой норме.
\end{proof}

\subsection{Якобиан}

Для отображения
\[
    g\colon\R^n\to\R^m,
    \qquad
    g(x)
    =
    \begin{pmatrix}
        g_1(x)\\
        \vdots\\
        g_m(x)
    \end{pmatrix},
\]
якобиан определяется как
\[
    J_g(x)
    =
    \frac{\partial g}{\partial x}
    =
    \begin{pmatrix}
        \frac{\partial g_1}{\partial x_1}
        &\cdots&
        \frac{\partial g_1}{\partial x_n}\\
        \vdots&&\vdots\\
        \frac{\partial g_m}{\partial x_1}
        &\cdots&
        \frac{\partial g_m}{\partial x_n}
    \end{pmatrix}
    \in\R^{m\times n}.
\]

\begin{theorem}[Многомерное правило цепочки]
Пусть $g\colon\R^n\to\R^m$ дифференцируема в точке $x$, а
$h\colon\R^m\to\R$ дифференцируема в точке $g(x)$. Тогда композиция
$F=h\circ g$ дифференцируема в точке $x$ и
\[
    \nabla F(x)
    =
    J_g(x)^\top
    \nabla h\bigl(g(x)\bigr).
\]
\end{theorem}

\begin{proof}
Обозначим $y=g(x)$. Из дифференцируемости $g$ следует
\[
    g(x+\Delta x)
    =
    g(x)
    +
    J_g(x)\Delta x
    +
    r_g(\Delta x),
\]
где
\[
    \frac{\norm{r_g(\Delta x)}}{\norm{\Delta x}}
    \longrightarrow 0
    \quad
    \text{при }
    \Delta x\to 0.
\]

Положим
\[
    \Delta y
    =
    g(x+\Delta x)-g(x)
    =
    J_g(x)\Delta x+r_g(\Delta x).
\]
Поскольку линейное отображение ограничено,
$\norm{\Delta y}=O(\norm{\Delta x})$.

Из дифференцируемости $h$ следует
\[
    h(y+\Delta y)
    =
    h(y)
    +
    \nabla h(y)^\top\Delta y
    +
    r_h(\Delta y),
\]
где
\[
    r_h(\Delta y)
    =
    o\bigl(\norm{\Delta y}\bigr)
    =
    o\bigl(\norm{\Delta x}\bigr).
\]

Подставляя выражение для $\Delta y$, получаем
\[
\begin{aligned}
    h(g(x+\Delta x))
    &=
    h(g(x))
    +
    \nabla h(g(x))^\top
    J_g(x)\Delta x\\
    &\quad+
    \nabla h(g(x))^\top r_g(\Delta x)
    +
    r_h(\Delta y).
\end{aligned}
\]
Два последних слагаемых имеют порядок
$o(\norm{\Delta x})$. Следовательно,
\[
    F(x+\Delta x)
    =
    F(x)
    +
    \left(
        J_g(x)^\top\nabla h(g(x))
    \right)^\top
    \Delta x
    +
    o(\norm{\Delta x}).
\]
По определению градиента
\[
    \nabla F(x)
    =
    J_g(x)^\top\nabla h(g(x)).
\]
\end{proof}

% ============================================================
\section{Градиентный спуск}
% ============================================================

\subsection{Основной алгоритм}

Пусть требуется минимизировать дифференцируемую функцию
\[
    F\colon\R^d\to\R.
\]
Метод градиентного спуска строит последовательность
\[
    \theta_{k+1}
    =
    \theta_k
    -
    \eta_k\nabla F(\theta_k),
\]
где $\eta_k>0$ --- шаг, или скорость обучения.

\begin{algorithm}[H]
\caption{Полный градиентный спуск}
\begin{algorithmic}[1]
\Require начальная точка $\theta_0$, число итераций $K$,
шаги $\eta_0,\ldots,\eta_{K-1}$
\Ensure приближённое решение $\theta_K$
\For{$k=0,\ldots,K-1$}
    \State вычислить значение $F(\theta_k)$
    \State вычислить полный градиент
    $g_k\gets\nabla F(\theta_k)$
    \If{выполнен критерий остановки}
        \State \Return $\theta_k$
    \EndIf
    \State обновить параметры:
    $\theta_{k+1}\gets\theta_k-\eta_k g_k$
\EndFor
\State \Return $\theta_K$
\end{algorithmic}
\end{algorithm}

В качестве критериев остановки могут использоваться:

\[
    \norm{\nabla F(\theta_k)}\leq\varepsilon,
\]
\[
    \abs{F(\theta_{k+1})-F(\theta_k)}
    \leq
    \varepsilon\max\{1,\abs{F(\theta_k)}\},
\]
ограничение на число итераций или отсутствие улучшения на
валидационной выборке.

\subsection{Минимальный одномерный пример}

Рассмотрим
\[
    F(\theta)
    =
    \frac{1}{2}(\theta-3)^2.
\]
Тогда
\[
    F'(\theta)=\theta-3
\]
и шаг градиентного спуска равен
\[
\begin{aligned}
    \theta_{k+1}
    &=
    \theta_k-\eta(\theta_k-3)\\
    &=
    (1-\eta)\theta_k+3\eta.
\end{aligned}
\]

При $\theta_0=0$ и $\eta=0.5$:
\[
\begin{aligned}
    \theta_1&=1.5,\\
    \theta_2&=2.25,\\
    \theta_3&=2.625,\\
    \theta_4&=2.8125.
\end{aligned}
\]

\begin{center}
\begin{tikzpicture}[x=1.7cm,y=1cm]
    \draw[-{Latex[length=2mm]}] (-0.25,0) -- (3.45,0)
        node[right] {$\theta$};

    \foreach \x/\lab in {
        0/{\theta_0},
        1.5/{\theta_1},
        2.25/{\theta_2},
        2.625/{\theta_3},
        3/{\theta_*}
    }{
        \draw (\x,0.08)--(\x,-0.08);
        \node[below] at (\x,-0.08) {$\lab$};
    }

    \draw[fwd] (0,0.20) to[bend left=25] (1.5,0.20);
    \draw[fwd] (1.5,0.20) to[bend left=25] (2.25,0.20);
    \draw[fwd] (2.25,0.20) to[bend left=25] (2.625,0.20);
\end{tikzpicture}
\end{center}

Ошибка относительно оптимума удовлетворяет равенству
\[
    \theta_{k+1}-3
    =
    (1-\eta)(\theta_k-3).
\]
Поэтому сходимость имеет место тогда, когда
\[
    \abs{1-\eta}<1,
    \qquad
    0<\eta<2.
\]

При слишком малом шаге сходимость медленная. При слишком большом
шаге итерации могут колебаться или расходиться.

\subsection{Гладкость функции}

\begin{definition}
Дифференцируемая функция $F\colon\R^d\to\R$ называется
$L$-гладкой, если её градиент является липшицевым:
\[
    \norm{\nabla F(x)-\nabla F(y)}_2
    \leq
    L\norm{x-y}_2
\]
для любых $x,y\in\R^d$.
\end{definition}

Константа $L$ ограничивает скорость изменения градиента. Большое
значение $L$ соответствует возможности высокой кривизны функции и
обычно требует меньшего шага градиентного спуска.

\begin{lemma}[Лемма о спуске]
Пусть $F\colon\R^d\to\R$ является $L$-гладкой. Тогда для любых
$x,y\in\R^d$
\[
    F(y)
    \leq
    F(x)
    +
    \nabla F(x)^\top(y-x)
    +
    \frac{L}{2}\norm{y-x}_2^2.
\]
\end{lemma}

\begin{proof}
Рассмотрим функцию одной переменной
\[
    \phi(t)
    =
    F\bigl(x+t(y-x)\bigr),
    \qquad t\in[0,1].
\]
По основной теореме анализа
\[
\begin{aligned}
    F(y)-F(x)
    &=
    \phi(1)-\phi(0)\\
    &=
    \int_0^1\phi'(t)\dd t\\
    &=
    \int_0^1
    \nabla F\bigl(x+t(y-x)\bigr)^\top
    (y-x)\dd t.
\end{aligned}
\]

Добавим и вычтем $\nabla F(x)$:
\[
\begin{aligned}
    F(y)-F(x)
    &=
    \nabla F(x)^\top(y-x)\\
    &\quad+
    \int_0^1
    \left[
        \nabla F(x+t(y-x))-\nabla F(x)
    \right]^\top
    (y-x)\dd t.
\end{aligned}
\]

По неравенству Коши--Буняковского и $L$-гладкости
\[
\begin{aligned}
    &
    \left[
        \nabla F(x+t(y-x))-\nabla F(x)
    \right]^\top
    (y-x)\\
    &\leq
    \norm{
        \nabla F(x+t(y-x))-\nabla F(x)
    }
    \norm{y-x}\\
    &\leq
    Lt\norm{y-x}^2.
\end{aligned}
\]

Следовательно,
\[
\begin{aligned}
    F(y)-F(x)
    &\leq
    \nabla F(x)^\top(y-x)
    +
    \int_0^1Lt\norm{y-x}^2\dd t\\
    &=
    \nabla F(x)^\top(y-x)
    +
    \frac{L}{2}\norm{y-x}^2.
\end{aligned}
\]
\end{proof}

\begin{corollary}[Убывание функции при одном шаге]
Если $F$ является $L$-гладкой и
\[
    x^+=x-\eta\nabla F(x),
\]
то
\[
    F(x^+)
    \leq
    F(x)
    -
    \eta
    \left(
        1-\frac{L\eta}{2}
    \right)
    \norm{\nabla F(x)}_2^2.
\]
В частности, если
\[
    0<\eta<\frac{2}{L},
\]
то при $\nabla F(x)\neq 0$ значение функции строго уменьшается.
\end{corollary}

\begin{proof}
В лемме о спуске положим
\[
    y=x-\eta\nabla F(x).
\]
Тогда
\[
    y-x=-\eta\nabla F(x),
\]
и поэтому
\[
\begin{aligned}
    F(y)
    &\leq
    F(x)
    -
    \eta\norm{\nabla F(x)}^2
    +
    \frac{L\eta^2}{2}
    \norm{\nabla F(x)}^2\\
    &=
    F(x)
    -
    \eta
    \left(
        1-\frac{L\eta}{2}
    \right)
    \norm{\nabla F(x)}^2.
\end{aligned}
\]
\end{proof}

\subsection{Сходимость в невыпуклом случае}

Функция потерь глубокой нейронной сети обычно невыпукла. Поэтому
градиентный спуск в общем случае не гарантирует нахождение
глобального минимума. Можно, однако, получить гарантию приближения
к стационарной точке.

\begin{theorem}[Сходимость градиента к нулю]
Пусть $F\colon\R^d\to\R$ является $L$-гладкой и ограничена снизу:
\[
    F_{\inf}
    =
    \inf_{x\in\R^d}F(x)>-\infty.
\]
Пусть последовательность определяется градиентным спуском
\[
    x_{k+1}=x_k-\eta\nabla F(x_k)
\]
с постоянным шагом
\[
    0<\eta<\frac{2}{L}.
\]
Тогда:

\begin{enumerate}
    \item последовательность $F(x_k)$ не возрастает;
    \item
    \[
        \sum_{k=0}^{\infty}
        \norm{\nabla F(x_k)}_2^2<\infty;
    \]
    \item
    \[
        \norm{\nabla F(x_k)}_2\longrightarrow 0;
    \]
    \item для любого $K\geq 1$
    \[
        \min_{0\leq k<K}
        \norm{\nabla F(x_k)}_2^2
        \leq
        \frac{
            F(x_0)-F_{\inf}
        }{
            \eta
            \left(
                1-\frac{L\eta}{2}
            \right)K
        }.
    \]
\end{enumerate}
\end{theorem}

\begin{proof}
По следствию из леммы о спуске
\[
    F(x_{k+1})
    \leq
    F(x_k)
    -
    c\norm{\nabla F(x_k)}^2,
\]
где
\[
    c
    =
    \eta
    \left(
        1-\frac{L\eta}{2}
    \right)>0.
\]
Отсюда
\[
    F(x_{k+1})\leq F(x_k),
\]
то есть последовательность значений функции не возрастает.

Перенесём слагаемые:
\[
    c\norm{\nabla F(x_k)}^2
    \leq
    F(x_k)-F(x_{k+1}).
\]
Суммируя от $k=0$ до $K-1$, получаем телескопическую сумму:
\[
\begin{aligned}
    c\sum_{k=0}^{K-1}
    \norm{\nabla F(x_k)}^2
    &\leq
    \sum_{k=0}^{K-1}
    \left(
        F(x_k)-F(x_{k+1})
    \right)\\
    &=
    F(x_0)-F(x_K)\\
    &\leq
    F(x_0)-F_{\inf}.
\end{aligned}
\]
Следовательно,
\[
    \sum_{k=0}^{K-1}
    \norm{\nabla F(x_k)}^2
    \leq
    \frac{F(x_0)-F_{\inf}}{c}.
\]
Правая часть не зависит от $K$. Переходя к пределу
$K\to\infty$, получаем
\[
    \sum_{k=0}^{\infty}
    \norm{\nabla F(x_k)}^2<\infty.
\]

Члены сходящегося ряда из неотрицательных чисел стремятся к нулю,
поэтому
\[
    \norm{\nabla F(x_k)}^2\to 0.
\]

Наконец,
\[
    K
    \min_{0\leq k<K}
    \norm{\nabla F(x_k)}^2
    \leq
    \sum_{k=0}^{K-1}
    \norm{\nabla F(x_k)}^2.
\]
Используя полученную оценку суммы, имеем
\[
    \min_{0\leq k<K}
    \norm{\nabla F(x_k)}^2
    \leq
    \frac{
        F(x_0)-F_{\inf}
    }{
        cK
    }.
\]
Подстановка значения $c$ завершает доказательство.
\end{proof}

\begin{remark}
Теорема утверждает сходимость нормы градиента к нулю, но не
гарантирует:

\begin{itemize}
    \item сходимость самих точек $x_k$;
    \item достижение глобального минимума;
    \item отличие найденной стационарной точки от седловой точки
    или локального максимума.
\end{itemize}
\end{remark}

\begin{remark}
Для сетей с ReLU функция потерь может быть недифференцируема на
границах линейных областей. Поэтому приведённая теорема в буквальном
виде относится к гладким функциям. На практике в точках излома
используется выбранное значение субградиента.
\end{remark}

\subsection{Скорость сходимости для выпуклой функции}

\begin{definition}
Функция $F\colon\R^d\to\R$ называется выпуклой, если для любых
$x,y\in\R^d$ и $\alpha\in[0,1]$
\[
    F(\alpha x+(1-\alpha)y)
    \leq
    \alpha F(x)+(1-\alpha)F(y).
\]
Если $F$ дифференцируема, выпуклость эквивалентна условию
\[
    F(y)
    \geq
    F(x)+\nabla F(x)^\top(y-x)
\]
для всех $x,y$.
\end{definition}

\begin{theorem}[Скорость градиентного спуска в выпуклом случае]
Пусть $F\colon\R^d\to\R$ выпукла, $L$-гладка и имеет точку
минимума $x_*$. Пусть
\[
    x_{k+1}
    =
    x_k-\frac{1}{L}\nabla F(x_k).
\]
Тогда для любого $k\geq 1$
\[
    F(x_k)-F(x_*)
    \leq
    \frac{
        L\norm{x_0-x_*}_2^2
    }{
        2k
    }.
\]
\end{theorem}

\begin{proof}
Обозначим
\[
    g_k=\nabla F(x_k).
\]
Из выпуклости следует
\[
    F(x_k)-F(x_*)
    \leq
    g_k^\top(x_k-x_*).
\]

Рассмотрим расстояние до минимума после одного шага:
\[
\begin{aligned}
    \norm{x_{k+1}-x_*}^2
    &=
    \norm{
        x_k-\frac{1}{L}g_k-x_*
    }^2\\
    &=
    \norm{x_k-x_*}^2
    -
    \frac{2}{L}
    g_k^\top(x_k-x_*)
    +
    \frac{1}{L^2}\norm{g_k}^2.
\end{aligned}
\]

Следовательно,
\[
\begin{aligned}
    \frac{L}{2}
    \left(
        \norm{x_k-x_*}^2
        -
        \norm{x_{k+1}-x_*}^2
    \right)
    &=
    g_k^\top(x_k-x_*)
    -
    \frac{1}{2L}\norm{g_k}^2\\
    &\geq
    F(x_k)-F(x_*)
    -
    \frac{1}{2L}\norm{g_k}^2.
\end{aligned}
\]

По лемме о спуске при шаге $1/L$
\[
    F(x_{k+1})
    \leq
    F(x_k)-\frac{1}{2L}\norm{g_k}^2.
\]
Поэтому
\[
    F(x_k)-F(x_*)
    -
    \frac{1}{2L}\norm{g_k}^2
    \geq
    F(x_{k+1})-F(x_*).
\]
Значит,
\[
    F(x_{k+1})-F(x_*)
    \leq
    \frac{L}{2}
    \left(
        \norm{x_k-x_*}^2
        -
        \norm{x_{k+1}-x_*}^2
    \right).
\]

Суммируем это неравенство по $k=0,\ldots,K-1$:
\[
\begin{aligned}
    \sum_{k=0}^{K-1}
    \left(
        F(x_{k+1})-F(x_*)
    \right)
    &\leq
    \frac{L}{2}
    \left(
        \norm{x_0-x_*}^2
        -
        \norm{x_K-x_*}^2
    \right)\\
    &\leq
    \frac{L}{2}\norm{x_0-x_*}^2.
\end{aligned}
\]

Значения функции не возрастают, поэтому
\[
    F(x_{k+1})-F(x_*)
    \geq
    F(x_K)-F(x_*)
\]
для всех $k<K$. Следовательно,
\[
    K\left(F(x_K)-F(x_*)\right)
    \leq
    \frac{L}{2}\norm{x_0-x_*}^2.
\]
Деление на $K$ завершает доказательство.
\end{proof}

\subsection{Стохастический и мини-пакетный градиентный спуск}

Пусть
\[
    F(\theta)
    =
    \frac{1}{N}
    \sum_{i=1}^{N}F_i(\theta).
\]

Полный градиент:
\[
    \nabla F(\theta)
    =
    \frac{1}{N}
    \sum_{i=1}^{N}\nabla F_i(\theta).
\]

Вычисление полного градиента может быть дорого при большом $N$.
В стохастическом градиентном спуске на каждой итерации выбирается
случайный объект $i_k$:
\[
    \theta_{k+1}
    =
    \theta_k
    -
    \eta_k\nabla F_{i_k}(\theta_k).
\]

Если $i_k$ равномерно распределён на
$\{1,\ldots,N\}$, то
\[
\begin{aligned}
    \E_{i_k}
    \left[
        \nabla F_{i_k}(\theta)
    \right]
    &=
    \sum_{i=1}^{N}
    \frac{1}{N}\nabla F_i(\theta)\\
    &=
    \nabla F(\theta).
\end{aligned}
\]
Таким образом, стохастический градиент является несмещённой оценкой
полного градиента.

На практике используется мини-пакет
\[
    B_k\subseteq\{1,\ldots,N\},
    \qquad
    \abs{B_k}=m,
\]
и оценка
\[
    g_{B_k}(\theta)
    =
    \frac{1}{m}
    \sum_{i\in B_k}
    \nabla F_i(\theta).
\]

Обновление:
\[
    \theta_{k+1}
    =
    \theta_k-\eta_k g_{B_k}(\theta_k).
\]

Если отдельные градиенты выбираются независимо и
\[
    \E\norm{\nabla F_i(\theta)-\nabla F(\theta)}^2
    =
    \sigma^2,
\]
то для среднего по $m$ независимым объектам
\[
    \E\norm{
        g_{B_k}(\theta)-\nabla F(\theta)
    }^2
    =
    \frac{\sigma^2}{m}.
\]
Увеличение размера мини-пакета уменьшает дисперсию оценки градиента,
но увеличивает стоимость одной итерации.

\begin{definition}
\textbf{Эпоха} --- обработка количества объектов, равного размеру
всей обучающей выборки. Если размер мини-пакета равен $m$, одна эпоха
содержит приблизительно $N/m$ итераций.
\end{definition}

\begin{algorithm}[H]
\caption{Обучение мини-пакетным градиентным спуском}
\begin{algorithmic}[1]
\Require данные $\mathcal{D}$, параметры $\theta$, число эпох $E$,
размер пакета $m$, скорости обучения $\eta_t$
\Ensure обученные параметры $\theta$
\For{$e=1,\ldots,E$}
    \State случайно перемешать обучающую выборку
    \State разбить выборку на мини-пакеты $B_1,\ldots,B_s$
    \For{$j=1,\ldots,s$}
        \State выполнить прямой проход для объектов из $B_j$
        \State вычислить среднюю функцию потерь
        \[
            F_{B_j}(\theta)
            =
            \frac{1}{\abs{B_j}}
            \sum_{i\in B_j}
            \ell(f_\theta(x_i),y_i)
        \]
        \State с помощью обратного распространения вычислить
        $g\gets\nabla_\theta F_{B_j}(\theta)$
        \State обновить параметры:
        $\theta\gets\theta-\eta_t g$
    \EndFor
    \State оценить функцию потерь и качество на валидационной выборке
\EndFor
\State \Return $\theta$
\end{algorithmic}
\end{algorithm}

\subsection{Выбор скорости обучения}

Основные стратегии:

\begin{itemize}
    \item постоянный шаг $\eta_k=\eta$;
    \item ступенчатое уменьшение шага;
    \item экспоненциальное уменьшение:
    \[
        \eta_k=\eta_0\gamma^k,
        \qquad 0<\gamma<1;
    \]
    \item полиномиальное уменьшение:
    \[
        \eta_k
        =
        \frac{\eta_0}{(1+\alpha k)^p};
    \]
    \item периодическое изменение;
    \item предварительный разогрев, при котором шаг постепенно
    увеличивается в начале обучения.
\end{itemize}

Признаки слишком большого шага:

\begin{itemize}
    \item функция потерь растёт или становится \texttt{NaN};
    \item нормы весов и градиентов быстро увеличиваются;
    \item значение функции сильно колеблется.
\end{itemize}

Признаки слишком малого шага:

\begin{itemize}
    \item функция потерь уменьшается крайне медленно;
    \item градиенты ненулевые, но параметры почти не меняются;
    \item обучение быстро выходит на длинное плато.
\end{itemize}

\subsection{Метод момента}

Метод момента сглаживает направления последовательных градиентов.

\begin{enumerate}
    \item Инициализировать $v_0=0$.
    \item На итерации $k$ вычислить градиент $g_k$.
    \item Обновить накопленный момент:
    \[
        v_{k+1}
        =
        \beta v_k+(1-\beta)g_k,
        \qquad 0\leq\beta<1.
    \]
    \item Обновить параметры:
    \[
        \theta_{k+1}
        =
        \theta_k-\eta v_{k+1}.
    \]
\end{enumerate}

Момент накапливает устойчивую составляющую градиента и подавляет
быстро меняющую знак составляющую.

\subsection{Adam}

Adam использует экспоненциальные средние градиента и квадрата
градиента. Все операции над векторами выполняются покомпонентно.

\begin{enumerate}
    \item Инициализировать
    \[
        m_0=0,
        \qquad
        v_0=0.
    \]
    \item На итерации $t$ вычислить $g_t=\nabla F_{B_t}(\theta_{t-1})$.
    \item Обновить первый момент:
    \[
        m_t
        =
        \beta_1m_{t-1}
        +
        (1-\beta_1)g_t.
    \]
    \item Обновить второй момент:
    \[
        v_t
        =
        \beta_2v_{t-1}
        +
        (1-\beta_2)g_t\odot g_t.
    \]
    \item Выполнить коррекцию начального смещения:
    \[
        \widehat{m}_t
        =
        \frac{m_t}{1-\beta_1^t},
        \qquad
        \widehat{v}_t
        =
        \frac{v_t}{1-\beta_2^t}.
    \]
    \item Обновить параметры:
    \[
        \theta_t
        =
        \theta_{t-1}
        -
        \eta
        \frac{\widehat{m}_t}
        {\sqrt{\widehat{v}_t}+\varepsilon}.
    \]
\end{enumerate}

Второй момент масштабирует шаг отдельно для каждой координаты.
Параметр $\varepsilon>0$ предотвращает деление на ноль.

% ============================================================
\section{Вычислительный граф}
% ============================================================

\subsection{Определение}

Вычислительный граф --- ориентированный ациклический граф, в котором:

\begin{itemize}
    \item входные вершины содержат исходные данные и параметры;
    \item внутренние вершины соответствуют элементарным операциям;
    \item рёбра описывают зависимости между результатами операций;
    \item выходная вершина содержит значение функции потерь.
\end{itemize}

Например,
\[
    u=wx+b,
    \qquad
    a=\varphi(u),
    \qquad
    \mathcal{L}=\frac{1}{2}(a-y)^2
\]
можно представить графом
\[
    (w,x)\longrightarrow wx
    \longrightarrow wx+b
    \longrightarrow \varphi(wx+b)
    \longrightarrow
    \frac{1}{2}(a-y)^2.
\]

\subsection{Прямой и обратный режимы автоматического
дифференцирования}

В \textbf{прямом режиме} вместе с каждым промежуточным значением
распространяется его производная по выбранному входу. Такой режим
удобен, когда входов мало, а выходов много.

В \textbf{обратном режиме} сначала вычисляются и сохраняются
промежуточные значения, после чего от скалярного выхода к входам
распространяются сопряжённые величины. Этот режим удобен, когда
выход скалярный, а параметров много. Именно такая ситуация
возникает при обучении нейронных сетей.

Для вершины $v$ определим сопряжённую величину
\[
    \overline{v}
    =
    \nabla_v\mathcal{L}.
\]

Если
\[
    v=g(u_1,\ldots,u_r),
\]
то вклад вершины $v$ в сопряжённую величину её родителя $u_j$ равен
\[
    \overline{u_j}
    \mathrel{+}=
    \left(
        \frac{\partial v}{\partial u_j}
    \right)^\top
    \overline{v}.
\]

Знак $\mathrel{+}=$ важен: если величина используется несколькими
последующими операциями, градиенты по всем путям складываются.

\begin{algorithm}[H]
\caption{Обратный режим дифференцирования вычислительного графа}
\begin{algorithmic}[1]
\Require вычислительный граф с параметрами $\theta$ и скалярным
выходом $\mathcal{L}$
\Ensure градиенты $\nabla_\theta\mathcal{L}$
\State упорядочить вершины топологически
\ForAll{вершин $v_i$ в прямом топологическом порядке}
    \State вычислить и сохранить значение $v_i$
\EndFor
\ForAll{вершин $v_i$}
    \State $\overline{v_i}\gets 0$
\EndFor
\State $\overline{\mathcal{L}}\gets 1$
\ForAll{вершин $v_i$ в обратном топологическом порядке}
    \ForAll{непосредственных родителей $u$ вершины $v_i$}
        \State
        \[
            \overline{u}
            \gets
            \overline{u}
            +
            \left(
                \frac{\partial v_i}{\partial u}
            \right)^\top
            \overline{v_i}
        \]
    \EndFor
\EndFor
\State \Return сопряжённые величины параметров
$\overline{\theta}=\nabla_\theta\mathcal{L}$
\end{algorithmic}
\end{algorithm}

\begin{theorem}[Корректность обратного накопления]
Пусть вычислительный граф является конечным ориентированным
ациклическим графом со скалярным выходом $\mathcal{L}$. Пусть все
локальные операции дифференцируемы в рассматриваемой точке. Тогда
алгоритм обратного режима вычисляет для каждой вершины $v$ точную
производную
\[
    \overline{v}
    =
    \nabla_v\mathcal{L}.
\]
\end{theorem}

\begin{proof}
Рассмотрим вершину $v$ и множество её непосредственных потомков
$\operatorname{Ch}(v)$. Любой ориентированный путь от $v$ к
$\mathcal{L}$ должен сначала пройти через одного из непосредственных
потомков $c\in\operatorname{Ch}(v)$.

По многомерному правилу цепочки полный градиент по $v$ равен сумме
вкладов всех непосредственных потомков:
\[
    \nabla_v\mathcal{L}
    =
    \sum_{c\in\operatorname{Ch}(v)}
    \left(
        \frac{\partial c}{\partial v}
    \right)^\top
    \nabla_c\mathcal{L}.
\]

Проведём индукцию по обратному топологическому порядку.

Для выходной вершины
\[
    \overline{\mathcal{L}}
    =
    \frac{\partial\mathcal{L}}
    {\partial\mathcal{L}}
    =
    1,
\]
поэтому база индукции выполнена.

Предположим, что к моменту обработки вершины $v$ для всех её
потомков уже вычислены точные величины
\[
    \overline{c}
    =
    \nabla_c\mathcal{L}.
\]
Это верно, поскольку в обратном топологическом порядке потомки
обрабатываются раньше родителей.

При обработке всех потомков в $\overline{v}$ последовательно
добавляются величины
\[
    \left(
        \frac{\partial c}{\partial v}
    \right)^\top
    \overline{c}.
\]
После учёта всех потомков
\[
\begin{aligned}
    \overline{v}
    &=
    \sum_{c\in\operatorname{Ch}(v)}
    \left(
        \frac{\partial c}{\partial v}
    \right)^\top
    \overline{c}\\
    &=
    \sum_{c\in\operatorname{Ch}(v)}
    \left(
        \frac{\partial c}{\partial v}
    \right)^\top
    \nabla_c\mathcal{L}\\
    &=
    \nabla_v\mathcal{L}.
\end{aligned}
\]
Следовательно, по индукции утверждение справедливо для всех вершин.
\end{proof}

% ============================================================
\section{Обратное распространение в многослойном перцептроне}
% ============================================================

\subsection{Обозначения}

Пусть
\[
    a^{(0)}=x,
\]
а для $l=1,\ldots,L$
\[
    z^{(l)}
    =
    W^{(l)}a^{(l-1)}+b^{(l)},
    \qquad
    a^{(l)}
    =
    \varphi_l(z^{(l)}).
\]

Пусть функция потерь для одного объекта равна
\[
    \mathcal{L}
    =
    \ell(a^{(L)},y).
\]

Определим ошибку слоя, или локальный градиент по предактивации:
\[
    \delta^{(l)}
    =
    \nabla_{z^{(l)}}\mathcal{L}.
\]

\subsection{Основные формулы}

\begin{theorem}[Формулы обратного распространения]
Предположим, что все функции активации и функция потерь
дифференцируемы в рассматриваемой точке. Тогда:

\begin{enumerate}
    \item для выходного слоя
    \[
        \delta^{(L)}
        =
        \nabla_{a^{(L)}}\ell
        \odot
        \varphi_L'\bigl(z^{(L)}\bigr);
    \]

    \item для скрытых слоёв
    \[
        \delta^{(l)}
        =
        \left(
            W^{(l+1)}
        \right)^\top
        \delta^{(l+1)}
        \odot
        \varphi_l'\bigl(z^{(l)}\bigr),
    \]
    где $l=L-1,\ldots,1$;

    \item градиенты по весам равны
    \[
        \nabla_{W^{(l)}}\mathcal{L}
        =
        \delta^{(l)}
        \left(
            a^{(l-1)}
        \right)^\top;
    \]

    \item градиенты по смещениям равны
    \[
        \nabla_{b^{(l)}}\mathcal{L}
        =
        \delta^{(l)}.
    \]
\end{enumerate}
\end{theorem}

\begin{proof}
Для покомпонентной функции активации введём диагональную матрицу
\[
    D^{(l)}
    =
    \diag\left(
        \varphi_l'(z_1^{(l)}),
        \ldots,
        \varphi_l'(z_{n_l}^{(l)})
    \right).
\]
Тогда дифференциал активации равен
\[
    \dd a^{(l)}
    =
    D^{(l)}\dd z^{(l)}.
\]

Начнём с выходного слоя. Дифференциал функции потерь:
\[
    \dd\mathcal{L}
    =
    \left(
        \nabla_{a^{(L)}}\ell
    \right)^\top
    \dd a^{(L)}.
\]
Так как
\[
    \dd a^{(L)}
    =
    D^{(L)}\dd z^{(L)},
\]
получаем
\[
\begin{aligned}
    \dd\mathcal{L}
    &=
    \left(
        \nabla_{a^{(L)}}\ell
    \right)^\top
    D^{(L)}
    \dd z^{(L)}\\
    &=
    \left(
        D^{(L)}
        \nabla_{a^{(L)}}\ell
    \right)^\top
    \dd z^{(L)}.
\end{aligned}
\]
Следовательно,
\[
    \delta^{(L)}
    =
    D^{(L)}
    \nabla_{a^{(L)}}\ell
    =
    \nabla_{a^{(L)}}\ell
    \odot
    \varphi_L'(z^{(L)}).
\]

Рассмотрим теперь слой $l+1$. Его предактивация:
\[
    z^{(l+1)}
    =
    W^{(l+1)}a^{(l)}+b^{(l+1)}.
\]
Полный дифференциал:
\[
\begin{aligned}
    \dd z^{(l+1)}
    &=
    \dd W^{(l+1)}a^{(l)}
    +
    W^{(l+1)}\dd a^{(l)}
    +
    \dd b^{(l+1)}.
\end{aligned}
\]

Часть дифференциала функции потерь, проходящая через этот слой:
\[
    \dd\mathcal{L}
    =
    \left(
        \delta^{(l+1)}
    \right)^\top
    \dd z^{(l+1)}.
\]
Подставим выражение для $\dd z^{(l+1)}$:
\[
\begin{aligned}
    \dd\mathcal{L}
    &=
    \left(
        \delta^{(l+1)}
    \right)^\top
    \dd W^{(l+1)}a^{(l)}\\
    &\quad+
    \left(
        \delta^{(l+1)}
    \right)^\top
    W^{(l+1)}\dd a^{(l)}\\
    &\quad+
    \left(
        \delta^{(l+1)}
    \right)^\top
    \dd b^{(l+1)}.
\end{aligned}
\]

Слагаемое, содержащее $\dd a^{(l)}$, преобразуется следующим
образом:
\[
\begin{aligned}
    \left(
        \delta^{(l+1)}
    \right)^\top
    W^{(l+1)}\dd a^{(l)}
    &=
    \left[
        \left(
            W^{(l+1)}
        \right)^\top
        \delta^{(l+1)}
    \right]^\top
    \dd a^{(l)}\\
    &=
    \left[
        D^{(l)}
        \left(
            W^{(l+1)}
        \right)^\top
        \delta^{(l+1)}
    \right]^\top
    \dd z^{(l)}.
\end{aligned}
\]

Следовательно,
\[
    \delta^{(l)}
    =
    D^{(l)}
    \left(
        W^{(l+1)}
    \right)^\top
    \delta^{(l+1)}.
\]
Поскольку $D^{(l)}$ диагональна, это эквивалентно
\[
    \delta^{(l)}
    =
    \left(
        W^{(l+1)}
    \right)^\top
    \delta^{(l+1)}
    \odot
    \varphi_l'(z^{(l)}).
\]

Осталось получить градиенты по параметрам. Для любого слоя
\[
    z^{(l)}
    =
    W^{(l)}a^{(l-1)}+b^{(l)}
\]
и
\[
    \dd\mathcal{L}
    =
    \left(
        \delta^{(l)}
    \right)^\top
    \dd z^{(l)}.
\]
Часть, зависящая от матрицы весов:
\[
    \left(
        \delta^{(l)}
    \right)^\top
    \dd W^{(l)}a^{(l-1)}.
\]
Используя скалярное произведение Фробениуса,
\[
    \left(
        \delta^{(l)}
    \right)^\top
    \dd W^{(l)}a^{(l-1)}
    =
    \inner{
        \delta^{(l)}
        (a^{(l-1)})^\top
    }{
        \dd W^{(l)}
    }_F.
\]
По определению градиента
\[
    \nabla_{W^{(l)}}\mathcal{L}
    =
    \delta^{(l)}
    (a^{(l-1)})^\top.
\]

Аналогично,
\[
    \left(
        \delta^{(l)}
    \right)^\top
    \dd b^{(l)}
\]
даёт
\[
    \nabla_{b^{(l)}}\mathcal{L}
    =
    \delta^{(l)}.
\]
\end{proof}

\subsection{Алгоритм обратного распространения}

\begin{algorithm}[H]
\caption{Обратное распространение для одного объекта}
\begin{algorithmic}[1]
\Require вход $x$, ответ $y$, параметры
$\{W^{(l)},b^{(l)}\}_{l=1}^{L}$
\Ensure значение потерь и градиенты по всем параметрам
\State $a^{(0)}\gets x$
\For{$l=1,\ldots,L$}
    \State
    $z^{(l)}
    \gets
    W^{(l)}a^{(l-1)}+b^{(l)}$
    \State
    $a^{(l)}
    \gets
    \varphi_l(z^{(l)})$
    \State сохранить $z^{(l)}$ и $a^{(l)}$
\EndFor
\State
$\mathcal{L}\gets\ell(a^{(L)},y)$
\State
\[
    \delta^{(L)}
    \gets
    \nabla_{a^{(L)}}\ell
    \odot
    \varphi_L'(z^{(L)})
\]
\For{$l=L,L-1,\ldots,1$}
    \State
    \[
        \nabla_{W^{(l)}}\mathcal{L}
        \gets
        \delta^{(l)}
        (a^{(l-1)})^\top
    \]
    \State
    \[
        \nabla_{b^{(l)}}\mathcal{L}
        \gets
        \delta^{(l)}
    \]
    \If{$l>1$}
        \State
        \[
            \delta^{(l-1)}
            \gets
            (W^{(l)})^\top\delta^{(l)}
            \odot
            \varphi_{l-1}'(z^{(l-1)})
        \]
    \EndIf
\EndFor
\State \Return $\mathcal{L}$ и все градиенты
\end{algorithmic}
\end{algorithm}

\subsection{Softmax и перекрёстная энтропия}

\begin{proposition}
Пусть
\[
    p_j
    =
    \frac{e^{z_j}}
    {\sum_{r=1}^{C}e^{z_r}},
\]
а $y\in\R^C$ удовлетворяет условиям
\[
    y_j\geq 0,
    \qquad
    \sum_{j=1}^{C}y_j=1.
\]
Для функции потерь
\[
    \ell(z,y)
    =
    -\sum_{j=1}^{C}y_j\log p_j
\]
градиент по логитам равен
\[
    \nabla_z\ell=p-y.
\]
\end{proposition}

\begin{proof}
Имеем
\[
    \log p_j
    =
    z_j
    -
    \log
    \left(
        \sum_{r=1}^{C}e^{z_r}
    \right).
\]
Поэтому
\[
\begin{aligned}
    \ell(z,y)
    &=
    -\sum_{j=1}^{C}y_jz_j
    +
    \sum_{j=1}^{C}y_j
    \log
    \left(
        \sum_{r=1}^{C}e^{z_r}
    \right)\\
    &=
    -y^\top z
    +
    \log
    \left(
        \sum_{r=1}^{C}e^{z_r}
    \right),
\end{aligned}
\]
так как $\sum_jy_j=1$.

Для координаты $z_k$
\[
\begin{aligned}
    \frac{\partial\ell}{\partial z_k}
    &=
    -y_k
    +
    \frac{e^{z_k}}
    {\sum_{r=1}^{C}e^{z_r}}\\
    &=
    -y_k+p_k.
\end{aligned}
\]
Следовательно,
\[
    \nabla_z\ell=p-y.
\]
\end{proof}

Таким образом, при использовании softmax и перекрёстной энтропии
не требуется явно строить полный якобиан softmax:
\[
    \delta^{(L)}=p-y.
\]

\subsection{Матричная форма для мини-пакета}

Пусть мини-пакет содержит $m$ объектов, записанных по столбцам:
\[
    A^{(0)}
    =
    X
    \in\R^{n_0\times m}.
\]
Тогда прямой проход:
\[
    Z^{(l)}
    =
    W^{(l)}A^{(l-1)}
    +
    b^{(l)}\mathbf{1}_m^\top,
\]
\[
    A^{(l)}
    =
    \varphi_l(Z^{(l)}).
\]

Обозначим матрицу ошибок слоя:
\[
    \Delta^{(l)}
    =
    \begin{pmatrix}
        \delta_1^{(l)}&
        \cdots&
        \delta_m^{(l)}
    \end{pmatrix}.
\]

Тогда
\[
    \Delta^{(l)}
    =
    \left(
        W^{(l+1)}
    \right)^\top
    \Delta^{(l+1)}
    \odot
    \varphi_l'(Z^{(l)}).
\]

Если функция потерь усредняется по пакету, то
\[
    \nabla_{W^{(l)}}F_B
    =
    \frac{1}{m}
    \Delta^{(l)}
    \left(
        A^{(l-1)}
    \right)^\top,
\]
\[
    \nabla_{b^{(l)}}F_B
    =
    \frac{1}{m}
    \Delta^{(l)}\mathbf{1}_m.
\]

Размерности:
\[
\begin{aligned}
    \Delta^{(l)}
    &\in\R^{n_l\times m},\\
    A^{(l-1)}
    &\in\R^{n_{l-1}\times m},\\
    \Delta^{(l)}(A^{(l-1)})^\top
    &\in\R^{n_l\times n_{l-1}}.
\end{aligned}
\]

\subsection{Численный пример одного шага обучения}

Рассмотрим сеть с одним скрытым ReLU-нейроном:
\[
    z^{(1)}=W^{(1)}x+b^{(1)},
    \qquad
    a^{(1)}=\operatorname{ReLU}(z^{(1)}),
\]
\[
    \widehat{y}
    =
    z^{(2)}
    =
    W^{(2)}a^{(1)}+b^{(2)}.
\]

Пусть
\[
    x=
    \begin{pmatrix}
        2\\
        1
    \end{pmatrix},
    \qquad
    y=1,
\]
\[
    W^{(1)}
    =
    \begin{pmatrix}
        0.5&-0.25
    \end{pmatrix},
    \qquad
    b^{(1)}=0,
\]
\[
    W^{(2)}=1.2,
    \qquad
    b^{(2)}=-0.1.
\]

Функция потерь:
\[
    \mathcal{L}
    =
    \frac{1}{2}(\widehat{y}-y)^2.
\]

\subsubsection*{Шаг 1. Прямое распространение}

\[
\begin{aligned}
    z^{(1)}
    &=
    0.5\cdot 2-0.25\cdot 1\\
    &=0.75,
\end{aligned}
\]
\[
    a^{(1)}
    =
    \operatorname{ReLU}(0.75)
    =
    0.75.
\]

Выход:
\[
\begin{aligned}
    \widehat{y}
    &=
    1.2\cdot 0.75-0.1\\
    &=0.8.
\end{aligned}
\]

Потери:
\[
    \mathcal{L}
    =
    \frac{1}{2}(0.8-1)^2
    =
    0.02.
\]

\subsubsection*{Шаг 2. Ошибка выходного слоя}

Так как выход линейный,
\[
    \delta^{(2)}
    =
    \frac{\partial\mathcal{L}}
    {\partial z^{(2)}}
    =
    \widehat{y}-y
    =
    -0.2.
\]

Градиенты параметров второго слоя:
\[
    \frac{\partial\mathcal{L}}{\partial W^{(2)}}
    =
    \delta^{(2)}a^{(1)}
    =
    -0.2\cdot 0.75
    =
    -0.15,
\]
\[
    \frac{\partial\mathcal{L}}{\partial b^{(2)}}
    =
    \delta^{(2)}
    =
    -0.2.
\]

\subsubsection*{Шаг 3. Ошибка скрытого слоя}

Поскольку $z^{(1)}=0.75>0$,
\[
    \operatorname{ReLU}'(z^{(1)})=1.
\]
Следовательно,
\[
\begin{aligned}
    \delta^{(1)}
    &=
    W^{(2)}\delta^{(2)}
    \operatorname{ReLU}'(z^{(1)})\\
    &=
    1.2\cdot(-0.2)\cdot 1\\
    &=
    -0.24.
\end{aligned}
\]

Градиенты первого слоя:
\[
\begin{aligned}
    \frac{\partial\mathcal{L}}{\partial W^{(1)}}
    &=
    \delta^{(1)}x^\top\\
    &=
    -0.24
    \begin{pmatrix}
        2&1
    \end{pmatrix}\\
    &=
    \begin{pmatrix}
        -0.48&-0.24
    \end{pmatrix},
\end{aligned}
\]
\[
    \frac{\partial\mathcal{L}}{\partial b^{(1)}}
    =
    -0.24.
\]

\subsubsection*{Шаг 4. Обновление параметров}

Пусть $\eta=0.1$. Тогда
\[
\begin{aligned}
    W_{\mathrm{new}}^{(1)}
    &=
    W^{(1)}
    -
    \eta
    \frac{\partial\mathcal{L}}{\partial W^{(1)}}\\
    &=
    \begin{pmatrix}
        0.5&-0.25
    \end{pmatrix}
    -
    0.1
    \begin{pmatrix}
        -0.48&-0.24
    \end{pmatrix}\\
    &=
    \begin{pmatrix}
        0.548&-0.226
    \end{pmatrix},
\end{aligned}
\]
\[
    b_{\mathrm{new}}^{(1)}
    =
    0-0.1(-0.24)
    =
    0.024,
\]
\[
    W_{\mathrm{new}}^{(2)}
    =
    1.2-0.1(-0.15)
    =
    1.215,
\]
\[
    b_{\mathrm{new}}^{(2)}
    =
    -0.1-0.1(-0.2)
    =
    -0.08.
\]

\subsubsection*{Шаг 5. Проверка результата}

Новая скрытая предактивация:
\[
\begin{aligned}
    z_{\mathrm{new}}^{(1)}
    &=
    0.548\cdot 2
    -
    0.226\cdot 1
    +
    0.024\\
    &=
    0.894.
\end{aligned}
\]

Новый выход:
\[
\begin{aligned}
    \widehat{y}_{\mathrm{new}}
    &=
    1.215\cdot 0.894-0.08\\
    &=
    1.00621.
\end{aligned}
\]

Новые потери:
\[
\begin{aligned}
    \mathcal{L}_{\mathrm{new}}
    &=
    \frac{1}{2}(1.00621-1)^2\\
    &\approx
    1.93\cdot 10^{-5}.
\end{aligned}
\]

В данном примере один шаг уменьшил потери с $0.02$ приблизительно
до $1.93\cdot 10^{-5}$.

% ============================================================
\section{Проблемы обучения глубоких сетей}
% ============================================================

\subsection{Исчезающие и взрывающиеся градиенты}

Пусть
\[
    D^{(l)}
    =
    \diag\left(
        \varphi_l'(z^{(l)})
    \right).
\]
Из формул обратного распространения следует
\[
\begin{aligned}
    \delta^{(l)}
    &=
    D^{(l)}
    (W^{(l+1)})^\top
    D^{(l+1)}
    (W^{(l+2)})^\top
    \cdots\\
    &\qquad\cdots
    D^{(L)}
    \nabla_{a^{(L)}}\ell.
\end{aligned}
\]

Для спектральной нормы
\[
\begin{aligned}
    \norm{\delta^{(l)}}
    &\leq
    \norm{\nabla_{a^{(L)}}\ell}
    \prod_{r=l}^{L}\norm{D^{(r)}}
    \prod_{r=l+1}^{L}\norm{W^{(r)}}.
\end{aligned}
\]

Если большинство множителей имеет норму меньше единицы, градиент
экспоненциально уменьшается при распространении к первым слоям.
Это называется \textbf{исчезновением градиента}.

Если произведение норм быстро растёт, возникает
\textbf{взрыв градиента}.

Для сигмоиды
\[
    0<\sigma'(z)\leq\frac{1}{4}.
\]
Поэтому многократное умножение на производные сигмоиды может быстро
уменьшать градиент, особенно если нейроны находятся в насыщении.

Основные способы борьбы:

\begin{itemize}
    \item подходящая инициализация весов;
    \item ReLU и её модификации;
    \item нормализация промежуточных представлений;
    \item остаточные связи;
    \item ограничение нормы градиента;
    \item архитектуры со специальными путями передачи информации.
\end{itemize}

\subsection{Инициализация весов}

Если все веса слоя инициализировать одинаково, нейроны будут
получать одинаковые градиенты и сохранять одинаковые параметры.
Поэтому веса инициализируются случайно.

Рассмотрим
\[
    z_j
    =
    \sum_{i=1}^{n_{\mathrm{in}}}W_{ji}a_i.
\]
Пусть $W_{ji}$ и $a_i$ независимы, имеют нулевые средние и одинаковые
вторые моменты. Тогда приближённо
\[
    \E[z_j^2]
    =
    n_{\mathrm{in}}
    \E[W_{ji}^2]
    \E[a_i^2].
\]

Чтобы масштаб активаций не увеличивался и не уменьшался от слоя к
слою, следует выбирать дисперсию весов обратно пропорционально
числу входов.

\subsubsection{Инициализация Xavier}

Для линейных слоёв и функций, работающих в приблизительно линейном
режиме, используется
\[
    \operatorname{Var}(W_{ji})
    =
    \frac{2}
    {n_{\mathrm{in}}+n_{\mathrm{out}}}.
\]

Нормальный вариант:
\[
    W_{ji}
    \sim
    \mathcal{N}
    \left(
        0,
        \frac{2}
        {n_{\mathrm{in}}+n_{\mathrm{out}}}
    \right).
\]

Равномерный вариант:
\[
    W_{ji}
    \sim
    U\left[
        -\sqrt{\frac{6}
        {n_{\mathrm{in}}+n_{\mathrm{out}}}},
        \sqrt{\frac{6}
        {n_{\mathrm{in}}+n_{\mathrm{out}}}}
    \right].
\]

\subsubsection{Инициализация He}

Для ReLU приблизительно половина симметрично распределённых
предактиваций обнуляется. Для сохранения второго момента используется
\[
    \operatorname{Var}(W_{ji})
    =
    \frac{2}{n_{\mathrm{in}}}.
\]

Нормальный вариант:
\[
    W_{ji}
    \sim
    \mathcal{N}
    \left(
        0,
        \frac{2}{n_{\mathrm{in}}}
    \right).
\]

Равномерный вариант:
\[
    W_{ji}
    \sim
    U\left[
        -\sqrt{\frac{6}{n_{\mathrm{in}}}},
        \sqrt{\frac{6}{n_{\mathrm{in}}}}
    \right].
\]

Смещения часто инициализируют нулями.

\subsection{Мёртвые ReLU}

Если для нейрона ReLU
\[
    z<0
\]
на всех поступающих объектах, то
\[
    \operatorname{ReLU}(z)=0,
    \qquad
    \operatorname{ReLU}'(z)=0.
\]
Градиент его входных весов становится нулевым, и нейрон может
перестать обучаться.

Возможные меры:

\begin{itemize}
    \item уменьшение скорости обучения;
    \item корректная инициализация;
    \item Leaky ReLU или параметрическая ReLU;
    \item нормализация активаций.
\end{itemize}

\subsection{Нормализация}

Для мини-пакета скалярных активаций
$u_1,\ldots,u_m$ batch normalization вычисляет
\[
    \mu_B
    =
    \frac{1}{m}
    \sum_{i=1}^{m}u_i,
\]
\[
    \sigma_B^2
    =
    \frac{1}{m}
    \sum_{i=1}^{m}
    (u_i-\mu_B)^2,
\]
\[
    \widehat{u}_i
    =
    \frac{u_i-\mu_B}
    {\sqrt{\sigma_B^2+\varepsilon}},
\]
\[
    y_i
    =
    \gamma\widehat{u}_i+\beta.
\]

Параметры $\gamma$ и $\beta$ обучаются вместе с остальной сетью.
На этапе применения модели используются накопленные оценки
статистик обучающей выборки.

Layer normalization использует аналогичную нормализацию, но
вычисляет статистики по признакам одного объекта, а не по объектам
мини-пакета.

\subsection{Остаточные связи}

В остаточном блоке
\[
    h^{(l+1)}
    =
    h^{(l)}
    +
    F_l(h^{(l)}).
\]
Якобиан блока:
\[
    \frac{\partial h^{(l+1)}}{\partial h^{(l)}}
    =
    I+J_{F_l}(h^{(l)}).
\]

При обратном распространении
\[
    \nabla_{h^{(l)}}\mathcal{L}
    =
    \left(
        I+J_{F_l}(h^{(l)})
    \right)^\top
    \nabla_{h^{(l+1)}}\mathcal{L}.
\]
Слагаемое с единичной матрицей создаёт прямой путь передачи
градиента между слоями.

\subsection{Ограничение нормы градиента}

При взрыве градиента можно заменить $g$ на
\[
    \widetilde{g}
    =
    g\cdot
    \min
    \left(
        1,
        \frac{\tau}{\norm{g}_2}
    \right),
\]
где $\tau>0$ --- максимально допустимая норма.

Если $\norm{g}\leq\tau$, градиент не изменяется. Если
$\norm{g}>\tau$, его направление сохраняется, а норма становится
равной $\tau$.

\subsection{Регуляризация}

\subsubsection{$L_2$-регуляризация}

Пусть
\[
    F_{\mathrm{reg}}(\theta)
    =
    F(\theta)
    +
    \frac{\lambda}{2}\norm{\theta}_2^2.
\]
Тогда
\[
    \nabla F_{\mathrm{reg}}(\theta)
    =
    \nabla F(\theta)+\lambda\theta.
\]

Шаг обычного градиентного спуска:
\[
\begin{aligned}
    \theta_{k+1}
    &=
    \theta_k
    -
    \eta
    \left(
        \nabla F(\theta_k)+\lambda\theta_k
    \right)\\
    &=
    (1-\eta\lambda)\theta_k
    -
    \eta\nabla F(\theta_k).
\end{aligned}
\]

\subsubsection{Dropout}

Во время обучения для активаций генерируется случайная маска
\[
    m_i\sim\operatorname{Bernoulli}(1-p),
\]
после чего используется
\[
    \widetilde{a}
    =
    \frac{m\odot a}{1-p}.
\]
Множитель $1/(1-p)$ сохраняет математическое ожидание активации:
\[
    \E[\widetilde{a}]=a.
\]
На этапе применения сети случайная маска не используется.

\subsubsection{Ранняя остановка}

Обучение останавливают, если качество на валидационной выборке не
улучшается заданное число эпох. При этом сохраняют параметры,
соответствующие лучшему валидационному результату.

\subsubsection{Другие методы}

\begin{itemize}
    \item увеличение и преобразование обучающих данных;
    \item ограничение нормы весов;
    \item добавление шума;
    \item уменьшение размера модели;
    \item усреднение предсказаний нескольких моделей.
\end{itemize}

% ============================================================
\section{Проверка и отладка обратного распространения}
% ============================================================

\subsection{Численная проверка градиента}

Для координаты $\theta_j$ центральная разность имеет вид
\[
    g_j^{\mathrm{num}}
    =
    \frac{
        F(\theta+\varepsilon e_j)
        -
        F(\theta-\varepsilon e_j)
    }{
        2\varepsilon
    },
\]
где $e_j$ --- $j$-й базисный вектор.

Аналитический градиент:
\[
    g_j^{\mathrm{analytic}}
    =
    \frac{\partial F}{\partial\theta_j}.
\]

Относительная ошибка:
\[
    r_j
    =
    \frac{
        \abs{
            g_j^{\mathrm{analytic}}
            -
            g_j^{\mathrm{num}}
        }
    }{
        \abs{g_j^{\mathrm{analytic}}}
        +
        \abs{g_j^{\mathrm{num}}}
        +
        \delta
    },
\]
где $\delta>0$ предотвращает деление на ноль.

\begin{algorithm}[H]
\caption{Проверка градиента конечными разностями}
\begin{algorithmic}[1]
\Require параметры $\theta$, функция $F$, аналитический градиент
$g$, малое число $\varepsilon$
\Ensure оценка корректности градиента
\State отключить dropout и другие случайные преобразования
\State зафиксировать небольшой мини-пакет
\State вычислить аналитический градиент
$g\gets\nabla_\theta F(\theta)$
\State выбрать несколько координат $j$
\ForAll{выбранных координат $j$}
    \State
    $\theta^+\gets\theta+\varepsilon e_j$
    \State
    $\theta^-\gets\theta-\varepsilon e_j$
    \State
    \[
        g_j^{\mathrm{num}}
        \gets
        \frac{
            F(\theta^+)-F(\theta^-)
        }{
            2\varepsilon
        }
    \]
    \State сравнить $g_j^{\mathrm{num}}$ и $g_j$
\EndFor
\end{algorithmic}
\end{algorithm}

Проверку желательно выполнять:

\begin{itemize}
    \item в арифметике повышенной точности;
    \item на маленькой сети и малом пакете;
    \item вдали от точек недифференцируемости ReLU;
    \item только для части координат, поскольку метод требует
    дополнительных прямых проходов.
\end{itemize}

\subsection{Проверка размерностей}

Для каждого слоя должны выполняться соотношения
\[
\begin{aligned}
    W^{(l)}
    &\in\R^{n_l\times n_{l-1}},\\
    a^{(l-1)}
    &\in\R^{n_{l-1}},\\
    z^{(l)},a^{(l)},\delta^{(l)}
    &\in\R^{n_l},\\
    \nabla_{W^{(l)}}\mathcal{L}
    &\in\R^{n_l\times n_{l-1}},\\
    \nabla_{b^{(l)}}\mathcal{L}
    &\in\R^{n_l}.
\end{aligned}
\]

Формула
\[
    \nabla_{W^{(l)}}\mathcal{L}
    =
    \delta^{(l)}(a^{(l-1)})^\top
\]
является внешним произведением векторов и автоматически имеет
размер матрицы весов.

\subsection{Полезные диагностические величины}

Во время обучения следует отслеживать:

\begin{itemize}
    \item потери на обучающей и валидационной выборках;
    \item норму градиента
    \[
        \norm{\nabla_\theta F};
    \]
    \item нормы параметров;
    \item долю нулевых активаций ReLU;
    \item распределение активаций и предактиваций;
    \item долю правильных ответов или другую целевую метрику;
    \item наличие значений \texttt{NaN} и бесконечностей.
\end{itemize}

Типичные признаки ошибок:

\begin{itemize}
    \item \textbf{потери не меняются}: слишком малый шаг, нулевые
    градиенты, отсоединённый вычислительный граф;

    \item \textbf{потери сразу становятся \texttt{NaN}}:
    слишком большой шаг, переполнение экспоненты, логарифм нуля,
    деление на ноль;

    \item \textbf{обучающая ошибка уменьшается, а валидационная
    растёт}: переобучение;

    \item \textbf{обе ошибки велики}: недостаточная мощность модели,
    плохая оптимизация, неверные признаки или ошибки реализации.
\end{itemize}

% ============================================================
\section{Вычислительная сложность}
% ============================================================

Для полносвязного слоя
\[
    z^{(l)}
    =
    W^{(l)}a^{(l-1)}+b^{(l)}
\]
стоимость матрично-векторного умножения равна
\[
    O(n_l n_{l-1}).
\]

Стоимость полного прямого прохода:
\[
    O\left(
        \sum_{l=1}^{L}
        n_l n_{l-1}
    \right).
\]

При обратном проходе выполняются умножения
\[
    (W^{(l)})^\top\delta^{(l)}
\]
и внешние произведения
\[
    \delta^{(l)}(a^{(l-1)})^\top.
\]
Их суммарная асимптотическая сложность имеет тот же порядок:
\[
    O\left(
        \sum_{l=1}^{L}
        n_l n_{l-1}
    \right).
\]

Обратный режим особенно эффективен потому, что за один обратный
проход вычисляет производные скалярной функции потерь по всем
параметрам сети.

Во время стандартного обратного распространения необходимо хранить
активации и предактивации прямого прохода. Требуемая память для
активаций одного объекта имеет порядок
\[
    O\left(
        \sum_{l=0}^{L}n_l
    \right),
\]
не считая памяти для параметров. Для мини-пакета размера $m$ объём
памяти для активаций увеличивается приблизительно в $m$ раз.

Если памяти недостаточно, часть промежуточных результатов можно не
хранить, а повторно вычислять во время обратного прохода. Это
уменьшает расход памяти ценой дополнительных вычислений.

% ============================================================
\section{Практический цикл обучения}
% ============================================================

Полный процесс обучения глубокой нейронной сети можно организовать
следующим образом.

\begin{enumerate}
    \item \textbf{Подготовить данные.}
    \begin{itemize}
        \item разделить данные на обучающую, валидационную и
        тестовую части;
        \item нормализовать числовые признаки;
        \item определить представление категориальных признаков;
        \item исключить утечку информации между частями выборки.
    \end{itemize}

    \item \textbf{Выбрать архитектуру.}
    \begin{itemize}
        \item определить число слоёв и их ширину;
        \item выбрать функции активации;
        \item определить выходной слой в соответствии с задачей.
    \end{itemize}

    \item \textbf{Выбрать функцию потерь.}
    \begin{itemize}
        \item среднеквадратичная ошибка для регрессии;
        \item бинарная перекрёстная энтропия для двух классов;
        \item softmax и перекрёстная энтропия для нескольких классов.
    \end{itemize}

    \item \textbf{Инициализировать параметры.}
    \begin{itemize}
        \item Xavier для линейных или близких к линейным активаций;
        \item He для ReLU;
        \item смещения обычно инициализировать нулями.
    \end{itemize}

    \item \textbf{Выбрать алгоритм оптимизации.}
    \begin{itemize}
        \item мини-пакетный SGD;
        \item SGD с моментом;
        \item Adam или другой адаптивный метод.
    \end{itemize}

    \item \textbf{На каждом мини-пакете выполнить:}
    \begin{enumerate}
        \item прямое распространение;
        \item вычисление потерь;
        \item обратное распространение;
        \item при необходимости ограничение нормы градиента;
        \item обновление параметров.
    \end{enumerate}

    \item \textbf{Контролировать обучение.}
    \begin{itemize}
        \item сравнивать обучающую и валидационную ошибки;
        \item отслеживать нормы градиентов;
        \item корректировать скорость обучения;
        \item применять раннюю остановку.
    \end{itemize}

    \item \textbf{После выбора модели} один раз оценить итоговое
    качество на тестовой выборке.
\end{enumerate}

% ============================================================
\section{Ключевые различия понятий}
% ============================================================

\begin{center}
\begin{tabularx}{\textwidth}{
    >{\raggedright\arraybackslash}p{0.25\textwidth}
    X
    X
}
\toprule
\textbf{Понятие} &
\textbf{Назначение} &
\textbf{Результат}\\
\midrule

Прямое распространение &
Последовательное вычисление слоёв сети &
Предсказание, промежуточные активации и значение потерь\\

Обратное распространение &
Применение правила цепочки в обратном порядке &
Градиенты функции потерь по параметрам\\

Градиентный спуск &
Изменение параметров в направлении уменьшения функции &
Новая точка пространства параметров\\

Автоматическое дифференцирование &
Систематическое дифференцирование программы, заданной
вычислительным графом &
Производные без ручного раскрытия полной символьной формулы\\
\bottomrule
\end{tabularx}
\end{center}

Обратное распространение не является самостоятельным методом
оптимизации. Оно только вычисляет градиент
\[
    \nabla_\theta F(\theta).
\]
Градиентный спуск использует этот градиент:
\[
    \theta_{k+1}
    =
    \theta_k-\eta_k\nabla_\theta F(\theta_k).
\]

% ============================================================
\section{Итоговые формулы}
% ============================================================

\subsection*{Прямой проход}

\[
    a^{(0)}=x,
\]
\[
    z^{(l)}
    =
    W^{(l)}a^{(l-1)}+b^{(l)},
\]
\[
    a^{(l)}
    =
    \varphi_l(z^{(l)}).
\]

\subsection*{Обратный проход}

\[
    \delta^{(L)}
    =
    \nabla_{a^{(L)}}\ell
    \odot
    \varphi_L'(z^{(L)}),
\]
\[
    \delta^{(l)}
    =
    (W^{(l+1)})^\top
    \delta^{(l+1)}
    \odot
    \varphi_l'(z^{(l)}).
\]

\subsection*{Градиенты параметров}

\[
    \nabla_{W^{(l)}}\mathcal{L}
    =
    \delta^{(l)}(a^{(l-1)})^\top,
\]
\[
    \nabla_{b^{(l)}}\mathcal{L}
    =
    \delta^{(l)}.
\]

\subsection*{Обновление градиентным спуском}

\[
    W^{(l)}
    \gets
    W^{(l)}
    -
    \eta
    \nabla_{W^{(l)}}F,
\]
\[
    b^{(l)}
    \gets
    b^{(l)}
    -
    \eta
    \nabla_{b^{(l)}}F.
\]

\subsection*{Softmax и перекрёстная энтропия}

\[
    p=\softmax(z),
    \qquad
    \ell=-\sum_jy_j\log p_j,
\]
\[
    \nabla_z\ell=p-y.
\]

\subsection*{Главная логическая цепочка обучения}

\[
\boxed{
\text{данные}
\longrightarrow
\text{прямой проход}
\longrightarrow
\text{потери}
\longrightarrow
\text{обратный проход}
\longrightarrow
\text{градиент}
\longrightarrow
\text{обновление параметров}
}
\]

\end{document}
